% ============================================================
% 203_R_Operator_En_ch.tex
% The recursion operator R in FFGFT
% Formal translation of the derivation from Doc. 180
% Johann Pascher — May 2026
% ============================================================

\section*{Abstract}

The recursion operator $\mathcal{R}$ of FFGFT formalises the
fractal scaling structure in mode space. This document translates
the constant $\xi_0 = 4/30000$, derived on physical grounds in Doc.~180,
into the formal language of the operator and checks whether
the fixed-point equation $\mathcal{R}(\Phi_*) = \Phi_*$ reproduces
this uniqueness. It does not: in no reading does the fixed-point
condition yield $\xi_0 = 4/30000$; the value remains fixed by the
three conditions from Doc.~180.

%------------------------------------------------------------
\section{Starting Point: The Physical Derivation (Doc.~180)}
%------------------------------------------------------------

Doc.~180 derives $\xi_0 = 4/30000$ from three independent conditions:

\begin{enumerate}
  \item \textbf{Dimensional contribution $10^{-4}$:}
    Four spacetime dimensions yield $\xi \sim (10^{-1})^4 = 10^{-4}$.
  \item \textbf{Geometric contribution $4/3$:}
    Tetrahedral symmetry: $4$ vertices / $3$ edges per vertex $= 4/3$.
  \item \textbf{Self-consistency:}
    The e-p-$\mu$ system yields $\kappa = 7$ as the only
    consistent recursion depth ($0.04\,\%$ accuracy).
\end{enumerate}

These three conditions fix $\xi_0$ without any free parameter.
The task of this document: translation into the language of $\mathcal{R}$.

%------------------------------------------------------------
\section{Definition of the Recursion Operator}
%------------------------------------------------------------

\subsection{Diagonal Form in Mode Space}

The recursion operator $\mathcal{R}$ acts diagonally on the Hilbert space
$\mathcal{H} = \bigoplus_{n,l,j}\mathbb{C}\cdot\mathbf{e}_{n,l,j}$:
\begin{equation}
  \mathcal{R}(\mathbf{e}_{n,l,j})
  = \lambda_{n,l,j} \cdot \mathbf{e}_{n,l,j}
  \label{eq:R-diagonal}
\end{equation}

The eigenvalues $\lambda_{n,l,j}$ are determined by the mass formula:
\begin{equation}
  \lambda_{n,l,j} = \xi_0^{p(n,l,j)-1} \cdot K_{\text{frak}}
  \label{eq:eigenwerte}
\end{equation}

with $K_{\text{frak}} \approx 0.986$ (fractal correction).

\subsection{Action on the Field Operator}

On the mode operator $\Phi = \sum_i m_i P_i$ (Doc.~202),
$\mathcal{R}$ does not act as a uniform scaling map
$\mathcal{R}(\Phi) = \xi_0 \cdot \Phi$. The time-mass duality
symmetry reads
\begin{equation}
  t \to \xi_0 \cdot t, \qquad m \to \xi_0^{-1} \cdot m,
  \qquad \Phi \to \Phi .
  \label{eq:tmd-symmetrie}
\end{equation}
With $\Phi = \sum_i m_i P_i$ and $m \to m/\xi_0$ one would get
$\Phi/\xi_0$, not $\xi_0\Phi$; at the same time
\eqref{eq:tmd-symmetrie} states $\Phi \to \Phi$, and the eigenvalues
\eqref{eq:eigenwerte} $\xi_0^{p-1}K_{\text{frak}}$ are mode-dependent,
hence not a uniform scaling by $\xi_0$. These three statements are
mutually inconsistent; a uniform scaling of $\Phi$ under $\mathcal{R}$
is therefore not determined.

%------------------------------------------------------------
\section{The Fixed-Point Equation}
%------------------------------------------------------------

\subsection{Definition of the Fixed Point}

The fixed point $\Phi_*$ is the operator that is invariant
under $\mathcal{R}$:
\begin{equation}
  \mathcal{R}(\Phi_*) = \Phi_*
  \label{eq:fixpunkt}
\end{equation}

A uniform scaling $\mathcal{R}(\Phi) = s\,\Phi$ with
$s = \xi_0$ or $s = 1/\xi_0$ would have only this fixed point:
\begin{equation}
  s \cdot \Phi_* = \Phi_*
  \quad \Rightarrow \quad
  (s - 1)\,\Phi_* = 0
\end{equation}

Since $s \neq 1$, it follows that $\Phi_* = 0$ --- a trivial fixed point.

\subsection{The Non-Trivial Fixed Point through Self-Consistency}

The physically relevant fixed point arises from the
self-consistency condition of the mass formula. The operator
$\Phi_* = \sum_i m_i^{(*)} P_i$ is a fixed point of $\mathcal{R}$
if:
\begin{equation}
  m_i^{(*)} = r_i \cdot \xi_0^{p_i} \cdot v
  \quad \text{and} \quad
  \mathcal{R}(m_i^{(*)}) = m_i^{(*)}
  \label{eq:selbstkonsistenz}
\end{equation}

The second condition means:
\begin{equation}
  \xi_0^{p_i - 1} \cdot K_{\text{frak}} \cdot m_i^{(*)}
  = m_i^{(*)}
  \quad \Rightarrow \quad
  \xi_0^{p_i - 1} \cdot K_{\text{frak}} = 1
  \label{eq:fixpunkt-bedingung}
\end{equation}

\subsection{Solving for \texorpdfstring{$\xi_0$}{xi0}}

From \eqref{eq:fixpunkt-bedingung} it follows for each mode $i$:
\begin{equation}
  \xi_0 = K_{\text{frak}}^{-1/(p_i-1)}
  \label{eq:xi0-aus-fixpunkt}
\end{equation}

For $\xi_0$ to be \textbf{universal} (mode-independent), this
equation would have to be satisfiable for all modes simultaneously.
With a universal $K_{\text{frak}} \approx 0.986$, however, it yields
neither a common value nor $\xi_0 = 4/30000$ (Step~3 below).

%------------------------------------------------------------
\section{Connection to the Three Conditions from Doc.~180}
%------------------------------------------------------------

\subsection{Step 1: Dimensional Contribution}

The recursion operator scales lengths by $\xi_0$. In $D=4$
spacetime dimensions, the 4D volume scales by $\xi_0^4$:
\begin{equation}
  \mathcal{R}^{(4D)}(V_4) = \xi_0^4 \cdot V_4
\end{equation}

For the mass operator, which scales with $\ell^{-1}$:
\begin{equation}
  \xi_0^{\text{dim}} = (10^{-1})^4 = 10^{-4}
\end{equation}

\subsection{Step 2: Geometric Factor}

The tetrahedral symmetry of the 4D torus supplies the geometric factor
from the topology of the faces:

\begin{center}
{\footnotesize
\adjustbox{max width=\textwidth}{%
\begin{tabular}{lcl}
\toprule
\textbf{Geometry} & \textbf{Value} & \textbf{Origin} \\
\midrule
Vertices of the tetrahedron & 4 & Point symmetry \\
Edges per vertex & 3 & Spatial embedding \\
Geometric factor $\xi_0^{\text{geo}}$ & $4/3$ & Ratio \\
\bottomrule
\end{tabular}%
}}
\end{center}

\subsection{Step 3: Self-Consistency}

The fixed-point condition \eqref{eq:xi0-aus-fixpunkt} with the lepton exponents
$p_e = 3/2$, $p_\mu = 1$, $p_\tau = 2/3$ serves as a consistency check.
For the electron ($p_e - 1 = 1/2$):
\begin{equation}
  \xi_0 = K_{\text{frak}}^{-1/(3/2 - 1)}
  = K_{\text{frak}}^{-2}
  \approx (0.986)^{-2} \approx 1.03
\end{equation}
This lies far from $\xi_0 = 4/30000 \approx 1.333 \cdot 10^{-4}$;
the reading $K_{\text{frak}}^2 = (0.986)^2 = 0.972$ does not reach
the value either. In no reading does the fixed-point condition yield
$\xi_0 = 4/30000$.

$\xi_0$ is therefore fixed solely by the three conditions from
Doc.~180, with the e-p-$\mu$ system: the exponent $\kappa = 7$ is the
only integer recursion depth that satisfies all three conditions
simultaneously.

%------------------------------------------------------------
\section{The Recursion Formula}
%------------------------------------------------------------

The recursion operator generates the mass hierarchy through repeated
application:
\begin{equation}
  \mathcal{R}^{(k)}(\mathbf{e}_{n,l,j})
  = \xi_0^{k(p_{n,l,j}-1)} \cdot K_{\text{frak}}^k
  \cdot \mathbf{e}_{n,l,j}
  \label{eq:rekursion-k}
\end{equation}

For $k = \kappa = 7$ (recursion depth from Doc.~180):
\begin{equation}
  \mathcal{R}^{(7)}(\mathbf{e}_e) = \xi_0^{7(3/2-1)} \cdot K_{\text{frak}}^7
  \cdot \mathbf{e}_e
  = \xi_0^{7/2} \cdot K_{\text{frak}}^7 \cdot \mathbf{e}_e
\end{equation}

The self-consistency condition requires this expression to be
consistent with the known electron mass. It does not fix the value
$\xi_0 = 4/30000$; that value is taken over from Doc.~180.

%------------------------------------------------------------
\section{Summary: R as a Formal Framework}
%------------------------------------------------------------

\begin{center}
{\footnotesize
\adjustbox{max width=\textwidth}{%
\begin{tabular}{p{4.5cm}p{6cm}}
\toprule
\textbf{Physical derivation} (Doc.~180) &
\textbf{Operator language} (Doc.~203) \\
\midrule
$10^{-4}$ from 4D spacetime &
  $\mathcal{R}^{(4D)}$ scales by $\xi_0^4$ \\[4pt]
$4/3$ from tetrahedral geometry &
  Geometric factor in $K_{\text{frak}}$ \\[4pt]
$\kappa = 7$ from the e-p-$\mu$ system &
  $\mathcal{R}^{(7)}(\Phi_*) = \Phi_*$ fixed-point condition \\[4pt]
$\xi_0 = 4/30000$ unique &
  Taken over; the fixed-point condition does not yield it \\
\bottomrule
\end{tabular}%
}
}
\end{center}

The recursion operator $\mathcal{R}$ is therefore not a new postulate,
but the formal representation of the geometry derived in Doc.~180.
The task from Doc.~202, Problem~1 (eigenvalue form of $\mathcal{R}$),
is thereby formulated; $\xi_0$ itself is not obtained from the
fixed-point condition but taken over from Doc.~180.

%------------------------------------------------------------
\section{Note: Apparent Contradiction from a Category Violation}
%------------------------------------------------------------

The fixed-point condition $\mathcal{R}(\Phi_*) = \Phi_*$ suggests
requiring $\lambda_i = 1$ for every mode $i$. Substituting
\eqref{eq:eigenwerte} yields:
\[
  K_{\text{frak}}(e) = \xi_0^{-1/2} \approx 86.6, \quad
  K_{\text{frak}}(\mu) = 1, \quad
  K_{\text{frak}}(\tau) = \xi_0^{1/3} \approx 0.051
\]

These values are all different --- an apparent contradiction
to a universal $K_{\text{frak}} \approx 0.986$.

\textbf{However:} this is \textbf{the same category violation}
as in the non-closure theorem (Doc.~193). Here, expressions from
different mode spaces are being compared --- just as
$r_e \cdot r_\mu$ is a meaningless quantity because $r_e$ and
$r_\mu$ are labels of different algebraic spaces,
the comparison $K_{\text{frak}}(e) = K_{\text{frak}}(\mu)$
is inadmissible.

The fixed-point condition $\mathcal{R}(\Phi_*) = \Phi_*$ is a
\textbf{global} condition on the field operator as a whole.
It is \textbf{not} equivalent to the requirement
$\lambda_i = 1$ for each $i$ separately, since the projectors
$P_i$ are orthogonal and do not enforce a common scaling.

$K_{\text{frak}} \approx 0.986$ is a single geometric
constant of the fractal torus (Doc.~180). Its value $1-100\,\xi_0$
is confirmed by the comparison with $\alpha$; the form of the fractal
correction factor is open.


%------------------------------------------------------------
\section*{Notation}
%------------------------------------------------------------

\begin{center}
{\footnotesize
\resizebox{\textwidth}{!}{\begin{tabular}{lp{9cm}}
\toprule
\textbf{Symbol} & \textbf{Meaning} \\
\midrule
$\mathcal{R}$ & Recursion operator in mode space \\
$\lambda_{n,l,j}$ & Eigenvalue of $\mathcal{R}$ for mode $(n,l,j)$ \\
$\lambda_{n,l,j} = \xi_0^{p-1}\cdot K_{\text{frak}}$ & Explicit eigenvalue formula \\
$K_{\text{frak}} \approx 0.986$ & Fractal correction constant (universal, Doc.~180) \\
$\kappa = 7$ & Recursion depth (from e-p-$\mu$ self-consistency) \\
$\Phi_* = \sum_i m_i P_i$ & Fixed-point state of the operator \\
$\mathcal{R}(\Phi_*) = \Phi_*$ & Fixed-point condition \\
$p(n,l,j)$ & Exponent function from torus geometry \\
$\Delta p_{\mu\to\tau} = 1/3$ & Step in the geom.\ sequence $p_{n+1}=(2/3)p_n$; $p_\mu\cdot(1-2/3)=1/3$ \\
$\xi_0 = 4/30000$ & Fundamental geometric constant (Doc.~180) \\
$P_i = |\mathbf{e}_i\rangle\langle\mathbf{e}_i|$ & Projector onto mode $i$ (orthogonal) \\
$m_i = r_i\cdot\xi_0^{p_i}\cdot v$ & FFGFT mass formula \\
\bottomrule
\end{tabular}}}
\end{center}

\medskip\noindent
\textbf{References:} Doc.~006 ($r_i$, $p_i$), Doc.~051 (spin structure),
Doc.~180 ($\xi_0$ derivation), Doc.~189 (stage structure),
Doc.~193 (non-closure theorem), Doc.~202 (overview).
